% TOP_PROBLEM: {"id": 1348, "title": "Woodin's GCH below Strongly Compact Cardinals", "category_id": 10, "category": {"id": 10, "name": "set_theory", "display_name": "Set Theory", "description": "Foundations of mathematics, infinite sets, and cardinality.", "slug": "set-theory", "order_index": 10, "created_at": "2026-07-31T15:26:25.671Z"}, "status": "open"}
% FIRST_POSTED: 2026-10-01
% ATTEMPT_STATUS: unresolved
% Imported from: unsolvedmath_additional_500.tex; UnsolvedMath ID 1348
% !TeX program = lualatex
% Compile twice: lualatex 1348.tex
% Self-contained: no external bibliography, images, or \input files.
% Numeric section labels and visible numbers are original UnsolvedMath IDs.
\documentclass[11pt,a4paper]{article}
\usepackage[margin=24mm,headheight=14pt]{geometry}
\usepackage{fontspec}
\setmainfont{Latin Modern Roman}
\setsansfont{Latin Modern Sans}
\setmonofont{Latin Modern Mono}
\usepackage{amsmath,amssymb,mathtools}
\usepackage{unicode-math}
\setmathfont{Latin Modern Math}
\usepackage{microtype}
\usepackage{enumitem,xurl,needspace}
\usepackage[dvipsnames]{xcolor}
\usepackage{hyperref}
\hypersetup{unicode=true,colorlinks=true,linkcolor=MidnightBlue,urlcolor=MidnightBlue,
 pdftitle={UnsolvedMath Problem 1348},pdfauthor={Source-annotated compilation},bookmarksnumbered=true}
\usepackage{bookmark,tocloft,titlesec,fancyhdr}
\setlength{\cftsecnumwidth}{4.2em}
\cftsetpnumwidth{2.5em}
\cftsetrmarg{4em}
\titleformat{\section}{\Large\bfseries\raggedright}{\thesection}{0.7em}{}
\pagestyle{fancy}\fancyhf{}
\fancyhead[L]{\small\sffamily Additional UnsolvedMath statements}
\fancyhead[R]{\small Original catalogue IDs}
\fancyfoot[C]{\thepage}
\setlength{\parindent}{0pt}
\setlength{\parskip}{5pt plus 1pt}
\setlength{\emergencystretch}{3em}
\setcounter{tocdepth}{1}\setcounter{secnumdepth}{1}
\setlist[itemize]{leftmargin=3em,itemsep=4pt,topsep=4pt,parsep=0pt}
\allowdisplaybreaks
\newcommand{\N}{\mathbb{N}}
\newcommand{\Z}{\mathbb{Z}}
\newcommand{\Q}{\mathbb{Q}}
\newcommand{\R}{\mathbb{R}}
\newcommand{\C}{\mathbb{C}}
\newcommand{\F}{\mathbb{F}}
\newcommand{\A}{\mathbb{A}}
\newcommand{\PP}{\mathbb{P}}
\newcommand{\E}{\mathbb{E}}
\newcommand{\eps}{\varepsilon}
\newcommand{\dd}{\,\mathrm d}
\newcommand{\sref}[2]{\hyperlink{source:#1:#2}{[S#2]}}
\newif\ifshowcatalogue
\showcataloguetrue
\newcommand{\cataloguescope}[1]{\ifshowcatalogue
 \paragraph{Statement recorded in the supplied source.}
 #1\fi}
\begin{document}
% UnsolvedMath ID 1348; source catalogue code ST-002

\setcounter{section}{1347}

\section{Global GCH from GCH Below a Strongly Compact Cardinal}\label{1348}

{\small\sffamily UnsolvedMath ID 1348 · Set Theory}

\subsection{Definitions and mathematical statement}

\paragraph{Shared notation.}Unless a section says otherwise, $\N=\{1,2,\ldots\}$,
$\Z$ is the ring of integers, $\Q,\R,\C$ are the rational, real, and complex
fields, $[N]=\{1,\ldots,\lfloor N\rfloor\}$, and $\log$ is the natural
logarithm. For real functions of a parameter tending to infinity,
$f\ll g$ and $f=O(g)$ mean $|f|\leq Cg$ eventually for some fixed $C$;
$f\gg g$ reverses this comparison; $f=o(g)$ means $f/g\to0$;
$f\sim g$ means $f/g\to1$; and $f\asymp g$ means both order bounds.
A subscript on $O$ or $\ll$ specifies allowed constant dependence.
The expression $n^{o(1)}$ in an upper bound means growth smaller than
$n^\epsilon$ for each fixed $\epsilon>0$. Local definitions govern any
exception, finite-field convention, or use of the same symbol for another
object. An asymptotic claim is not automatically an effective algorithm.



For an infinite cardinal $\lambda$, GCH at $\lambda$ is $2^\lambda=\lambda^+$. An uncountable cardinal $\kappa$ is strongly compact when every $\kappa$-complete filter on any set extends to a $\kappa$-complete ultrafilter. Completeness means closure under intersections of fewer than $\kappa$ sets. The question asks whether GCH at every infinite $\lambda<\kappa$ forces GCH at every infinite cardinal. \sref{1348}{1} \sref{1348}{2}

\paragraph{Statement recorded in the supplied source.}
Does the generalized continuum hypothesis below a strongly compact cardinal imply it everywhere? \sref{1348}{1}

\subsection{Short English statement}

Can the generalized continuum hypothesis below one strongly compact cardinal force it to hold at all larger cardinals as well?

\subsection{Research attempt: locating the first possible GCH failure}
\textbf{Status and inputs.} The propagation to all cardinals is not proved. The inspected 2025 primary slides [S2] explicitly ask this Woodin question. We use one established external theorem: Solovay's singular-cardinal theorem says that above a strongly compact cardinal every singular strong-limit cardinal $\lambda$ satisfies $2^\lambda=\lambda^+$. A primary research account restating and reproving that theorem is [R1]. The cardinal-arithmetic reductions below are proved directly.

\subsubsection{1. GCH below a singular first failure makes it strong limit}
Suppose $\kappa$ is strongly compact, GCH holds below $\kappa$, and GCH fails somewhere. Let $\lambda$ be the least infinite cardinal with $2^\lambda>\lambda^+$. Such a least failure exists by the well-ordering of the cardinals, and necessarily $\lambda\geq\kappa$. If $\lambda$ were singular, it would be a limit cardinal. For every infinite $\mu<\lambda$, minimality would give
\[2^\mu=\mu^+<\lambda.\]
Thus $\lambda$ would be strong limit. Solovay's theorem then forbids the alleged failure. Consequently any least failure must be \emph{regular}. This conclusion invokes the large-cardinal theorem exactly once, at the singular strong-limit step. It does not invoke an unproved extension of that theorem to regular cardinals.

\subsubsection{2. The arithmetic identity behind the singular restriction}
For additional precision, let $\lambda$ be any infinite singular strong-limit cardinal and $\theta=\operatorname{cf}(\lambda)$. Choose a cofinal increasing sequence $(\lambda_i)_{i<\theta}$ of infinite cardinals below $\lambda$. A subset $A\subseteq\lambda$ is determined by $(A\cap\lambda_i)_{i<\theta}$, so
\[2^\lambda\leq\prod_{i<\theta}2^{\lambda_i}\leq\lambda^\theta.
\]
Conversely
\[\lambda^\theta\leq(2^\lambda)^\theta
=2^{\lambda\cdot\theta}=2^\lambda,
\]
using $\theta\leq\lambda$ and infinite cardinal multiplication in ZFC. Therefore
\[2^\lambda=\lambda^{\operatorname{cf}(\lambda)}.\tag{1}\]
This identifies the singular obstruction with a precise exponent, rather than with an informal extrapolation of the regular continuum function. It also shows why control of singular-cardinal arithmetic is substantial without supplying control at every regular cardinal.

\subsubsection{3. Why limit continuity is not a proof}
At such a strong-limit $\lambda$, one has $2^{<\lambda}=\sup_{\mu<\lambda}2^\mu=\lambda$. Cantor's theorem still gives $2^\lambda>\lambda$. Therefore the tempting identity $2^\lambda=\sup_{\mu<\lambda}2^\mu$ is false even in the favorable strong-limit setting. It confuses bounded subsets with arbitrary subsets coded by an entire cofinal sequence, whose number is (1).

Similarly, a least failure being regular does not contradict monotonicity of exponentiation: a jump above $\lambda^+$ can remain larger at all subsequent cardinals. The usual lower-cardinal equations by themselves provide no upper bound $2^\lambda\leq\lambda^+$ at that first regular point. Strong compactness would need to supply further structure specifically ruling out that jump.

\subsubsection{4. Verification and exact remaining statement}
The proof audit checks that the use of $\mu^+<\lambda$ is restricted to limit $\lambda$, that (1) uses the strong-limit hypothesis for every factor, and that the named external theorem applies only at singular strong limits. No finite computation can test these uncountable claims; their verification here is symbolic and conditional on the explicitly identified Solovay theorem. The unresolved reduced question is: can there be a regular $\lambda\geq\kappa$ that is the first failure of GCH while $\kappa$ is strongly compact and GCH holds below it? Eliminating such a regular first failure would complete the original implication; no elimination is given here.

\subsection{Sources}

{\small

\begin{itemize}

\item[\hypertarget{source:1348:1}{S1}] ULAM.AI, \emph{UnsolvedMath}, supplied \texttt{problems.json}, record 1348 (ST-002), statement and background. The embedded literature-review date is 2026-08-17; that is the archive’s review date, not a fresh certification. Dataset landing page: \url{https://huggingface.co/datasets/ulamai/UnsolvedMath}.

\item[\hypertarget{source:1348:2}{S2}] M. Golshani, Set theory questions (2025 lecture slides). \url{https://www.dmg.tuwien.ac.at/fb8/2025_slides/Golshani.pdf}. Bibliographic information imported from the supplied record.

\item[R1] Y. Matsubara, Saturated ideals and the singular cardinal hypothesis, Journal of Symbolic Logic; primary paper whose abstract states Solovay’s theorem and describes its generic-ultrapower proof.
\url{https://www.cambridge.org/core/services/aop-cambridge-core/content/view/9AA9D2B987612AA9A0574052FF5B057E/S0022481200022465a.pdf/saturated-ideals-and-the-singular-cardinal-hypothesis.pdf}.
\end{itemize}

}

\label{end:1348}
\end{document}
